\documentclass[11pt,a4paper]{article}
\usepackage[margin=25mm]{geometry}
\usepackage{amsmath,amssymb,amsthm,hyperref}
\setlength{\emergencystretch}{2em}
\newtheorem{theorem}{Theorem}
\newtheorem{lemma}[theorem]{Lemma}
\newcommand{\gen}[1]{\langle#1\rangle}
\newcommand{\F}{\mathbf F}
\newcommand{\Tr}{\operatorname{Tr}}
\title{Source interfaces for characteristic-two orthogonal markings}
\author{Human Compression project: L031 corrected recovery}
\date{October 2026}
\begin{document}\maketitle
This paper supplies the source-specific interfaces and auxiliary
qualifications for the complete corrected recovery of
Brooksbank--Ferrara--Leemans [BFL]. The full companion papers are
required: \emph{Endpoint variation in orthogonal symmetry markings},
\emph{The four displayed symplectic generators}, the two binary
boundary/exclusion papers, the centralizer rank-bound paper, and the
previously accepted \emph{Projective involutions and support spaces
in small dimension}. None is replaced by a reference to its conclusions
alone. The classical results identified below remain external inputs.

\section{Forms and the exact support boundary}
We use the source's definition of an orthogonal space: the polar
radical has dimension at most one. For a one-dimensional radical
$\gen z$ choose a nondegenerate complement $M$. If $Q(z)=0$, every
isometry is uniquely of the form
\[
 z\longmapsto cz,\qquad x\longmapsto hx+f(x)z\quad(x\in M),
 \qquad c\in F^\times, h\in O(M), f\in M^*.
\]
This proves the split description (2.3), with additive radical
$M^*$, identified with $M$ through the nondegenerate pairing.
If $Q(z)\ne0$, an isometry fixes $z$. Conversely each symplectic
map $h$ on $M$ has a unique lift
$x\mapsto hx+f(x)z$, where
$f(x)^2Q(z)=Q(x)+Q(hx)$. The right side is a quadratic function
with zero polar form; over the perfect field it is the square of a
unique linear function times $Q(z)$. This proves (2.4) directly.

In a nondegenerate plane, hyperbolic versus asingular means that
the binary quadratic polynomial has respectively two or no
projective roots. If the polar form is zero, the quadratic form
is a squared linear function and has either one singular point
or is identically zero. This gives all four source line types.
For nonradical nonsingular $x$, the formula
$s_x(v)=v+B(v,x)x/Q(x)$ is a nonidentity involution with support
$\gen x$. Distinct such symmetries commute exactly when their
root vectors are perpendicular, as follows by multiplying their
rank-one differences from the identity. Scaling the root by a
\emph{nonzero} scalar leaves the symmetry unchanged. A nonsingular
\emph{radical} root instead gives the identity; the source's generic
involution/support language must be read with this qualification.

\begin{theorem}[The larger-field intersection statement]
For $q\geq4$, the conclusion of [BFL, Proposition 3.3] holds in
every orthogonal ambient space allowed there, including a singular
one-dimensional radical.
\end{theorem}
\begin{proof}
For a nonsingular ambient space this is exactly the support
characterization in the endpoint companion, applied to the two
subspaces and their intersection. If its radical $\gen z$ is singular,
adjoin a vector $w$ with $B(z,w)=1$ and $Q(w)=0$, perpendicular
to a chosen nondegenerate complement of $z$. The larger space
is nondegenerate. All its symmetries rooted in the old space
preserve that space. Restriction of their whole group is faithful:
an element supported in the old space and fixing it pointwise has
the form $x\mapsto x+cB(x,z)z$; quadratic preservation and
$Q(z)=0$ force $c=0$. The proved ambient intersection equality
therefore restricts to the asserted one.
\end{proof}

\begin{theorem}[The unrestricted binary proposition is false]
The same Proposition 3.3, literally stated without excluding $q=2$,
has a counterexample with $\dim V=8$, $\dim W_i=6$ and
$\dim(W_1\cap W_2)=5$.
\end{theorem}
\begin{proof}
Let $M=M_2(\F_2)$ with $Q(X)=\det X$, the plus quadratic
four-space. Its nine nonzero singular vectors are the rank-one
matrices $uv^{\mathsf T}$, a $3$ by $3$ grid indexed by the
nonzero vectors $u,v\in\F_2^2$. Its maximal totally singular
planes are the three row fibers and three column fibers of this
grid: if two independent rank-one matrices have independent
left and right factors, their sum has rank two. Thus an isometry
permutes these two families, possibly interchanging them.
The product of the signs of its two induced $S_3$ permutations is
a well-defined homomorphism $\chi:O(M)\to\{1,-1\}$, also
when the families are interchanged.

Every nonsingular matrix $A$ has determinant one. The polar
identity for the determinant gives the symmetry formula
\[
                       s_A(X)=A\operatorname{adj}(X)A.
\]
Left and right multiplication by $A$ have the same ruling sign.
Moreover $\operatorname{adj}(X)=JX^{\mathsf T}J$, where
$J=\left(\begin{smallmatrix}0&1\\1&0\end{smallmatrix}\right)$;
the two $J$ signs cancel and transposition only swaps the families.
Consequently $\chi(s_A)=1$ for every symmetry.
Left multiplication $g:X\mapsto JX$, however, has $\chi(g)=-1$.
In particular $g\notin\Sigma_M$.

Now take $V=M\perp H_1\perp H_2$, where $H_i$ have singular
hyperbolic bases $e_i,f_i$ with $B(e_i,f_i)=1$. Put
\[
 U=M+\gen{e_1},\quad W_1=M+\gen{e_1,f_1},\quad
 W_2=M+\gen{e_1,f_1+e_2}.
\]
Both $W_i$ are nondegenerate six-spaces and their intersection is
$U$, whose radical is the singular line $\gen{e_1}$.
Extend $g$ by the identity on $H_1\perp H_2$. It is supported
in $M$ and preserves both $W_i$ and their perpendicular complements.
Each full six-dimensional orthogonal group is generated by
symmetries; this is the explicitly retained classical generation
result [Taylor, Corollary 11.42], whose sole exception is
$O^+(4,2)$. Therefore $g\in\Sigma_{W_1}\cap\Sigma_{W_2}$.
Every symmetry rooted in $U$ restricts on $U/\gen{e_1}=M$
to a symmetry of $M$. Since $g$ induces the already excluded
map on this quotient, $g\notin\Sigma_U$. This proves the
literal failure. No finite computation is a premise.
\end{proof}
The positive construction uses only the proved $q\geq4$ statement.
The two separate binary papers give the corrected existential
dimension boundary; they do not use Proposition 3.3.

\section{Three-spaces and full-order interval generation}
The intrinsic statement [BFL, Lemma 3.4] follows as well.
For nonsingular radical, $\Sigma_V=\mathrm{SL}_2(q)$ by (2.4),
and the asingular plane gives its maximal $D_{2(q+1)}$.
For singular radical, writing $V=L\perp\gen z$ identifies
$\Sigma_V$ with $L^+\rtimes O^-(L)$. All translations follow
from $s_v s_{v+az}$ as in the companion. The cyclic group of
order $q+1$ is irreducible on $L$ even over $\F_2$; hence any
subgroup properly containing the complement contains all translations.
This proves maximality for arbitrary elements, not only a chosen
third symmetry. When $q=2$ and the radical is nonsingular, the
two groups are both $S_3$, as stated in Remark 3.5.

For a hyperbolic plane and $q\geq8$, replace the torus by
$\operatorname{diag}(a,a^{-1})$, $a$ primitive in $\F_q^\times$,
and include the interchange of its two axes. The minimal
polynomials of $a$ and $a^{-1}$ over $\F_2$ are different:
equality would give $2^j\equiv-1\pmod{2^k-1}$ for some
$0\leq j<k$, impossible when $k\geq3$. They each have degree
$k$, so the two torus constituents are irreducible, nonisomorphic,
and interchanged by the involution. The full dihedral module is
therefore irreducible over $\F_2$. The singular-radical argument
just given works again. For nonsingular radical retain the split
torus-normalizer maximality in $\mathrm{SL}_2(q)$.
This proves the $q\geq8$ hyperbolic assertion in Remark 3.5.
For $q=4$, take $a\in\F_4$ of order three. The orbit of $(1,1)$
under $\gen{\operatorname{diag}(a,a^{-1}),\text{swap}}$ consists
of the three nonzero vectors of a proper two-dimensional
$\F_2$-subspace. Its translations with the complement form the
proper intermediate group asserted in that remark.

The rank-one maximality input here is the classical Dickson
subgroup classification: in even characteristic both full torus
normalizers $D_{2(q\pm1)}$ are maximal for $q\geq4$.
An exact locally available formulation is Proposition 2.3, page 5,
of Leemans--Moerenhout--O'Reilly-Regueiro,
\emph{Projective linear groups as automorphism groups of chiral
polytopes}, citing Bray--Holt--Roney-Dougal and Dickson.
It is retained classification, not proved by these constructions.

Now take the source path with every adjacent coefficient giving
order $q+1$. The endpoint companion proves generation of every
interval of length at least three; length two is its full dihedral
plane and length one is immediate. If the full ambient path has
singular radical, adjoin one more nonzero end coefficient to make
an even nondegenerate path. The same interval proofs hold in that
ambient space, and their restriction to the original space is
faithful by the preceding embedding argument. This proves all of
Lemma 5.3, and also its natural length-$d$ extension used in the
source proof. The full IP induction follows the same way; it
identifies the generated \emph{symmetry group} when the radical
is singular, not the additional radical scalings in (2.3).

For $q\geq8$ the mixed $q-1,q+1$ assertion of Remark 5.4 follows
by the same induction using the hyperbolic version just proved.
At each step use the first marked two-space $L$ and the suffix
hyperplane $W$. Their intersection contains the nonsingular
second basis root. Thus the intermediate two-space has at most
two singular points; exclude also the single point in $L$ and
the possible radical point to obtain the third symmetry needed
in the ambient three-space argument. In the nonsingular-radical
case its central kernel is recovered by the same Dickson-character
and perfect-$\mathrm{SL}_2$ argument as in the companion.
This proves the asserted generation without silently extending
an intrinsic three-space calculation to the ambient group.

For $q=4$, a bond of order three forces its normalized coefficient
to be one. A length-$b$ all-three path is therefore the scalar
extension of the binary path, whose group is $S_{b+1}$ by the
paid binary companion. Thus the source's symmetric-group
subscripts are one too small: $t+1$ consecutive generators give
$S_{t+2}$, and $d$ binary path generators give $S_{d+1}$.

\section{Radicals and the Arf prescription}
For the path $Q=\sum x_i^2+\sum\alpha_i x_ix_{i+1}$,
successively solving its tridiagonal polar equations proves
Lemma 4.2 in its full stated range. The even determinant is
$(\alpha_1\alpha_3\cdots\alpha_{d-1})^2$. In odd dimension
$2m+1$ the radical vector has alternate entries
\[
 z=(1,0,\beta_1,0,\ldots,0,\beta_m),\qquad
 \beta_j=\prod_{i=1}^j\alpha_{2i-1}/\alpha_{2i},
\]
and $Q(z)=1+\sum\beta_i^2$.

In even dimension set $e_1=v_1$ and, recursively,
\[
 e_i=v_{2i-1}+(\alpha_{2i-2}/\alpha_{2i-3})e_{i-1},
 \qquad f_i=v_{2i}/\alpha_{2i-1}.
\]
The odd basis vectors are mutually perpendicular, as are the
even ones. Direct cancellation gives $B(e_i,f_j)=\delta_{ij}$.
Expanding the recursion is exactly (4.5); quadratic evaluation
gives $Q(e_i)=1+(\alpha_{2i-2}/\alpha_{2i-3})^2Q(e_{i-1})$.
Thus the Arf trace is $\Tr\sum_i Q(e_i)/\alpha_{2i-1}^2$.
For $m\geq2$, equal endpoint coefficients $\lambda$ and equal
interior coefficients $\mu$, both full-order parameters, this is
\[
 \binom{m-1}{2}\Tr(\mu^{-1})+\Tr(\mu/\lambda^2).
\]
Since $\Tr(\lambda^{-1})=\Tr(\mu^{-1})=1$, this equals the
trace of the printed expression (5.6). Passing between inverse
squares and inverses is legitimate modulo $\{x^2+x:x\in F\}$.
The primitive norm-one eigenvalues, paired with their inverses,
give exactly $\varphi(q+1)/2$ full-order parameters, so their
existence and the source parameter sets are explicit.

The subsequent source instruction to choose a primitive $\mu$
with an additional trace constraint is not used as an existence
premise here. The complete endpoint-variation proof replaces that
instruction by a surjective trace calculation with freely variable
nonzero endpoints. Both Witt types and every stated dimension
are obtained. In dimension two the two dihedral groups are treated
directly, rather than substituting into a formula with an unused
interior coefficient. This is a corrected proof of the full
principal type-existence conclusion, not a claim that the source's
particular two-primitive-parameter prescription has been certified.

\section{Complete low-dimensional principal spectrum}
The accepted whole paper \emph{Projective involutions and support
spaces in small dimension}, from the L011 recovery, supplies the
even-characteristic exclusions for $\mathrm{PSL}_3(q)$,
$\mathrm{PSU}_3(q)$, $\mathrm{PSL}_4(q)$ and $\mathrm{PSU}_4(q)$.
Its canonical projective lifts and both six-dimensional quadratic
types are explicit. It retains the classical natural-module facts,
Dickson invariant and Klein identifications; none of its exclusions
uses the L031 theorem being recovered. The entire proof paper is
an acquired dependency here, not a circular citation to L011's
classification application.

For $\mathrm{SL}_2(q)$ in characteristic two, the centralizer of
a nonidentity involution is the abelian upper-unipotent group in
an appropriate basis. In a connected string of length four, its
first generator centralizes the noncommuting last pair, impossible.
This also excludes the natural rank-one classical identifications
appearing in dimensions three and four, including
$\Omega^-_4(q)\cong\mathrm{SL}_2(q^2)$.
The plus four-dimensional derived orthogonal group is a product
of two rank-one factors and is not quasisimple.

For the remaining quasisimple group $\operatorname{Sp}(4,q)$,
$q\geq4$, the endpoint companion gives rank five, the literal
four-matrix companion gives rank four, and the complete root-authored
centralizer companion excludes every connected string of length
six or more. Simplicity ensures connectedness for a full-group
string C-marking. These prove both directions of Theorem 6.1
for its stated ranks $r\geq4$, with no retained maximal-subgroup
rank assertion and no finite certificate.

For completeness its small-field aside also has an elementary
check. The full $\operatorname{Sp}(4,2)\cong S_6$ is not
quasisimple. In its derived group $A_6$, every involution is a
double transposition and has centralizer $D_8$. Thus a string
C-marking of rank at least five would place its last three of
the first five generators in a dihedral centralizer. An irreducible
rank-three C-marking cannot generate a dihedral subgroup: its
commuting endpoints differ by a central half-turn, contradicting
the rank-three intersection property. The rank-four case is equally
direct. Normalize the first generator to $(12)(34)$. A distinct
commuting last generator is, up to its centralizer,
either $(13)(24)$ or $(12)(56)$. In the first case the two middle
generators, required not to commute with their neighbours, belong
respectively to
\[
 \{(13)(56),(24)(56)\},\quad
 \{(12)(56),(34)(56)\}.
\]
All four marks preserve $\{1,2,3,4\}\mid\{5,6\}$.
In the second case their respective possibilities are
\[
 \{(15)(26),(16)(25)\},\quad
 \{(13)(24),(14)(23)\};
\]
all four preserve the three pairs $\{1,2\},\{3,4\},\{5,6\}$.
Neither case generates $A_6$. The lists follow simply by commuting
a double transposition with each normalized endpoint, and exhaust
the four other involutions of each $D_8$ centralizer.

\paragraph{What is not a recovered question.}
The odd-prime modular-reduction Theorem 1.3 is attributed prior work
and is not a new proof or a premise of the characteristic-two
construction. The regular-polytope/string-C-group correspondence
and the end-parabolic criterion are retained from
McMullen--Schulte. The historical availability of Magma code is
recorded; no such original executable is supplied or priced as zero.
The questions in Section 7 about other binary markings, higher
derived orthogonal groups and arbitrary-rank projective special
linear constructions remain source-recorded questions, with no
new answer or present-day open-status assertion. The corrected
binary theorem, auxiliary counterexample and intermediate plus
type are source corrections, not conjecture-credit claims.

\begin{thebibliography}{9}
\bibitem{BFL} P. A. Brooksbank, J. T. Ferrara and D. Leemans,
\emph{Orthogonal groups in characteristic 2 acting on polytopes
of high rank}, arXiv:1709.02219v2, supplied twelve-page edition.
\bibitem{Taylor} D. E. Taylor, \emph{The Geometry of the Classical
Groups}, Heldermann, 1992: Theorem 11.39 (larger-field symmetry
generation), Corollary 11.42 (the binary exception), Theorem 11.43
(Dickson invariant). Author-hosted primary representation:
\url{https://mail.maths.usyd.edu.au/u/don/papers/gcg.pdf},
printed pp.157--160, physical pp.165--168.
\bibitem{LMO} D. Leemans, J. Moerenhout and E. O'Reilly-Regueiro,
\emph{Projective linear groups as automorphism groups of chiral
polytopes}, supplied twenty-five-page edition, Proposition 2.3,
p.5; its rank-one classification is explicitly borrowed from
Bray--Holt--Roney-Dougal and Dickson, with external proof cost retained.
\bibitem{MS} P. McMullen and E. Schulte, \emph{Abstract Regular
Polytopes}, Cambridge, 2002, Section 2E, in particular Proposition 2E16.
\end{thebibliography}
\end{document}
